\documentclass[a4paper,11pt]{article} 
\usepackage[french]{babel} 
\usepackage{exptech} 
\usepackage{graphicx} 
\usepackage{booktabs}
\usepackage{array} 
\usepackage{hyperref}
\usepackage{listings} 
\usepackage{xcolor}
\usepackage{url}
\lstset{
    basicstyle=\ttfamily\small,
    breaklines=true,
    frame=single,
    backgroundcolor=\color{gray!10}
}
\title{
\begin{center}
{\Huge \textbf{Manuel Utilisateur}} \\[1cm]
{\textbf{} Aide à la lecture pour les manuels techniques d'appareils électroniques}
\end{center}
}
\author{Victoire Peltier, Ilyas El maaroufi, Maïwenn Lepéroux, Rayane El bouhali, Louis Richard
  \\ \\
  Encadrant : Jean-Louis Pazat}
\date{2025-2026}
\begin{document}
\maketitle
\section{Manuel utilisateur}
\subsection{Prérequis et Installation}
L'application est pleinement fonctionnelle sur \textbf{Linux}. Des instabilités peuvent survenir sur \textbf{macOS} (blocages lors de la fermeture de la fenêtre graphique) ; il est alors nécessaire de \textit{forcer à quitter} le processus Python.
Le projet est hébergé et accessible à l'adresse suivante : \url{https://gitlab.insa-rennes.fr/vpeltier/etudes-pratiques.git}.
Pour utiliser l'application, Tesseract OCR doit être installé sur votre système. Vous pouvez vous référer au guide officiel d'installation disponible à l'adresse : \url{https://tesseract-ocr.github.io/tessdoc/Installation.html}.
Ensuite, installez l'ensemble des dépendances Python requises en exécutant la commande suivante dans votre terminal :
\begin{lstlisting}[language=bash]
pip install ipywidgets pdf2image numpy opencv-python tqdm ultralytics pytesseract matplotlib easyocr
\end{lstlisting}
\subsection{Lancement et Utilisation}
\subsubsection{Via le script Python}
Il est possible de lancer l'application directement depuis le terminal, sans passer par JupyterLab, en utilisant le script \texttt{code\_complet.py} situé dans le dossier \texttt{Application\_finale}. Pour ce faire, exécutez la commande suivante :
\begin{lstlisting}[language=bash]
python3 Application_finale/code_complet.py <document.pdf>
\end{lstlisting}
où \texttt{<document.pdf>} est le nom du fichier PDF à traiter. 

\subsubsection{Via JupyterLab}
Pour démarrer l'application via JupyterLab, suivez les étapes suivantes :
\begin{enumerate}
    \item Ouvrez un terminal et accédez au dossier du projet :
    \begin{lstlisting}[language=bash]
cd Application_finale
    \end{lstlisting}
    \item Lancez l'environnement JupyterLab :
    \begin{lstlisting}[language=bash]
jupyter-lab
    \end{lstlisting}
    \item Dans JupyterLab, ouvrez le fichier \texttt{documentation\_technique.ipynb}.
    \item Modifiez si nécessaire le chemin du PDF à analyser ainsi que le dossier d'enregistrement. L'emplacement de ce code pour changer le PDF et le dossier d'enregistrement se trouve dans la toute dernière cellule du notebook.
    \item Exécutez toutes les cellules dans l'ordre. L'interface interactive s'affichera également au niveau de cette dernière cellule.
\end{enumerate}
\paragraph{Prise en main de l'interface :}
Une fois l'interface chargée, vous pouvez naviguer ainsi :
\begin{itemize}
    \item \textbf{Par la liste :} Cliquez sur un composant dans la liste de gauche pour recentrer automatiquement les trois vues sur celui-ci.
    \item \textbf{Dézoom :} Effectuez un clic droit pour dézoomer toutes les vues.
\end{itemize}
\end{document}

